\documentclass[profile=editorial]{atelier}
\AtelierUseACGNBibliography
\addbibresource{examples/acgn-research-demo/references.bib}

\title{ACGN 跨媒介研究示例}
\subtitle{游戏、动画、漫画、音乐与 Live 的统一引用}
\author{AtelierTeX Demo}

\begin{document}
\maketitle

\section{为什么需要媒体级定位}

二次元研究的材料往往不是单一“书/论文”。同一论证可能同时引用游戏本体\cite{demo_game}、具体角色剧情\cite{demo_commu}、动画单集\cite[00:12:31--00:12:46]{demo_anime}、漫画单话\cite[45--52]{demo_manga}、角色歌\cite{demo_song}与舞台演出\cite{demo_live}。

AtelierTeX 将正式 BibLaTeX entry type 与 `verba` 媒体导航标签分开，使书目规范和 ACGN 可读性可以同时保留。

\begin{AtelierOriginal}
日本語の一次資料は、必要な短い範囲だけ引用する。
\end{AtelierOriginal}

\begin{AtelierTranslation}
日文 primary text 只引用论证所需的短片段；长场景应通过精确 locator 转述与分析。
\end{AtelierTranslation}

\section{FAN LOCATOR 的位置}

玩家整理可以帮助找到原作位置\cite{demo_locator}，但当论文判断依赖作品事实时，应回到游戏、动画、漫画或其他 primary source 复核。

\printbibliography
\end{document}
